% GENERATED FILE. Edit canonical lessons, then run npm run generate:books.
% canonical-source-hash: fnv1a64:6ac9e4d775fa29a8
% canonical-lessons: SA-C493-samskarah, SA-C493-namakaranam, SA-C493-annaprasanam, SA-C493-cudakaranam, SA-C493-upanayanam

\chapter{Rites of Childhood}
\label{ch:sa-rites-of-childhood}
\hlchaptermodality{493}

\begin{chapteropening}
\textbf{By the end of this chapter:} \emph{I can name the rites of a child's early years.}

Complete the last lesson of chapter 493: I can name the rites of a child's early years.
\end{chapteropening}

\section[\texorpdfstring{\sk{संस्कारः} (saṃskāraḥ)}{संस्कारः (saṃskāraḥ)}]{\sk{संस्कारः} (saṃskāraḥ) — a rite of passage}

\label{lesson:SA-C493-samskarah}



\begin{quote}
Before the new one: say the Sanskrit for newborn, then the Sanskrit for a swing.
\end{quote}

\subsection*{You'll want to know: \sk{संस्कारः}}
\textbf{\sk{संस्कारः}} — \emph{saṃskāraḥ} — \textquotedblleft{}a rite of passage\textquotedblright{}.

\begin{quote}
\sk{विवाहः} \sk{अपि} \sk{एकः} \sk{संस्कारः} \sk{अस्ति।}
\end{quote}

\begin{quote}
\emph{vivāhaḥ api ekaḥ saṃskāraḥ asti.}
\end{quote}

\begin{quote}
A wedding too is a rite of passage.
\end{quote}

A masculine noun.

\begin{cousinweb}
\textbf{\sk{संस्कारः}} (\emph{saṃskāraḥ}) is built from \emph{sam} and the root √kṛ, \textquotedblleft{}to make\textquotedblright{}: a making perfect; MW gives \textquotedblleft{}a sacred or sanctifying ceremony\textquotedblright{}, twelve of them in a life (MW).
\end{cousinweb}

\begin{grammarlens}[title={What you've built}]
One more word for describing what you lived and felt.
\end{grammarlens}

\subsection*{Your turn}
\begin{itemize}
\raggedright
  \item \emph{Say it:} \emph{saṃskāraḥ}
  \item \emph{Say it:} \emph{saṃskāraḥ}, once more
  \item \emph{Say it:} \emph{vivāhaḥ api ekaḥ saṃskāraḥ asti.} — A wedding too is a rite of passage.
  \item \emph{Recall:} say the Sanskrit for newborn, then the Sanskrit for a swing, then say \emph{saṃskāraḥ} again
\end{itemize}

\begin{tcolorbox}[breakable,colback=teal!4,colframe=teal!35!black,title={Before you move on}]
What does \sk{संस्कारः} mean? (\textquotedblleft{}A rite of passage\textquotedblright{}.) Say it once more.
\end{tcolorbox}
\section[\texorpdfstring{\sk{नामकरणम्} (nāmakaraṇam)}{नामकरणम् (nāmakaraṇam)}]{\sk{नामकरणम्} (nāmakaraṇam) — the naming ceremony}

\label{lesson:SA-C493-namakaranam}



\begin{quote}
Before the new one: say the Sanskrit for a swing, then the Sanskrit for a rite of passage.
\end{quote}

\subsection*{You'll want to know: \sk{नामकरणम्}}
\textbf{\sk{नामकरणम्}} — \emph{nāmakaraṇam} — \textquotedblleft{}the naming ceremony\textquotedblright{}.

\begin{quote}
\sk{दशमे} \sk{दिने} \sk{शिशोः} \sk{नामकरणं} \sk{भवति।}
\end{quote}

\begin{quote}
\emph{daśame dine śiśoḥ nāmakaraṇaṃ bhavati.}
\end{quote}

\begin{quote}
On the tenth day the baby is given a name.
\end{quote}

A neuter noun.

\begin{cousinweb}
\textbf{\sk{नामकरणम्}} (\emph{nāmakaraṇam}) is built from \textbf{\sk{नाम}}, \textquotedblleft{}a name\textquotedblright{}, and \emph{karaṇa}, \textquotedblleft{}making\textquotedblright{}, from the root √kṛ; MW gives \textquotedblleft{}the ceremony of naming a child after birth\textquotedblright{} (MW).
\end{cousinweb}

\begin{grammarlens}[title={What you've built}]
One more word for describing what you lived and felt.
\end{grammarlens}

\subsection*{Your turn}
\begin{itemize}
\raggedright
  \item \emph{Say it:} \emph{nāmakaraṇam}
  \item \emph{Say it:} \emph{nāmakaraṇam}, once more
  \item \emph{Say it:} \emph{daśame dine śiśoḥ nāmakaraṇaṃ bhavati.} — On the tenth day the baby is given a name.
  \item \emph{Recall:} say the Sanskrit for a swing, then the Sanskrit for a rite of passage, then say \emph{nāmakaraṇam} again
\end{itemize}

\begin{tcolorbox}[breakable,colback=teal!4,colframe=teal!35!black,title={Before you move on}]
What does \sk{नामकरणम्} mean? (\textquotedblleft{}The naming ceremony\textquotedblright{}.) Say it once more.
\end{tcolorbox}
\section[\texorpdfstring{\sk{अन्नप्राशनम्} (annaprāśanam)}{अन्नप्राशनम् (annaprāśanam)}]{\sk{अन्नप्राशनम्} (annaprāśanam) — a baby's first rice}

\label{lesson:SA-C493-annaprasanam}



\begin{quote}
Before the new one: say the Sanskrit for a rite of passage, then the Sanskrit for the naming ceremony.
\end{quote}

\subsection*{You'll want to know: \sk{अन्नप्राशनम्}}
\textbf{\sk{अन्नप्राशनम्}} — \emph{annaprāśanam} — \textquotedblleft{}a baby's first rice\textquotedblright{}.

\begin{quote}
\sk{षष्ठे} \sk{मासे} \sk{शिशोः} \sk{अन्नप्राशनं} \sk{भवति।}
\end{quote}

\begin{quote}
\emph{ṣaṣṭhe māse śiśoḥ annaprāśanaṃ bhavati.}
\end{quote}

\begin{quote}
In the sixth month the baby has its first rice.
\end{quote}

A neuter noun.

\begin{cousinweb}
\textbf{\sk{अन्नप्राशनम्}} (\emph{annaprāśanam}) is built from \textbf{\sk{अन्नम्}}, \textquotedblleft{}food\textquotedblright{}, and \emph{prāśana}, \textquotedblleft{}feeding\textquotedblright{}, from \emph{pra} and the root √aś, \textquotedblleft{}to eat\textquotedblright{}; MW gives \textquotedblleft{}putting rice into a child's mouth for the first time\textquotedblright{} (MW).
\end{cousinweb}

\begin{grammarlens}[title={What you've built}]
One more word for describing what you lived and felt.
\end{grammarlens}

\subsection*{Your turn}
\begin{itemize}
\raggedright
  \item \emph{Say it:} \emph{annaprāśanam}
  \item \emph{Say it:} \emph{annaprāśanam}, once more
  \item \emph{Say it:} \emph{ṣaṣṭhe māse śiśoḥ annaprāśanaṃ bhavati.} — In the sixth month the baby has its first rice.
  \item \emph{Recall:} say the Sanskrit for a rite of passage, then the Sanskrit for the naming ceremony, then say \emph{annaprāśanam} again
\end{itemize}

\begin{tcolorbox}[breakable,colback=teal!4,colframe=teal!35!black,title={Before you move on}]
What does \sk{अन्नप्राशनम्} mean? (\textquotedblleft{}A baby's first rice\textquotedblright{}.) Say it once more.
\end{tcolorbox}
\section[\texorpdfstring{\sk{चूडाकरणम्} (cūḍākaraṇam)}{चूडाकरणम् (cūḍākaraṇam)}]{\sk{चूडाकरणम्} (cūḍākaraṇam) — the first haircut}

\label{lesson:SA-C493-cudakaranam}



\begin{quote}
Before the new one: say the Sanskrit for the naming ceremony, then the Sanskrit for a baby's first rice.
\end{quote}

\subsection*{You'll want to know: \sk{चूडाकरणम्}}
\textbf{\sk{चूडाकरणम्}} — \emph{cūḍākaraṇam} — \textquotedblleft{}the first haircut\textquotedblright{}.

\begin{quote}
\sk{चूडाकरणे} \sk{नापितः} \sk{शिशोः} \sk{केशान्} \sk{छिनत्ति।}
\end{quote}

\begin{quote}
\emph{cūḍākaraṇe nāpitaḥ śiśoḥ keśān chinatti.}
\end{quote}

\begin{quote}
At the first haircut the barber cuts the baby's hair.
\end{quote}

A neuter noun.

\begin{cousinweb}
\textbf{\sk{चूडाकरणम्}} (\emph{cūḍākaraṇam}) is built from \emph{cūḍā}, \textquotedblleft{}a tuft of hair\textquotedblright{}, and \emph{karaṇa}, \textquotedblleft{}making\textquotedblright{}; MW gives \textquotedblleft{}forming the crest, the ceremony of tonsure\textquotedblright{} (MW).
\end{cousinweb}

\begin{grammarlens}[title={What you've built}]
One more word for describing what you lived and felt.
\end{grammarlens}

\subsection*{Your turn}
\begin{itemize}
\raggedright
  \item \emph{Say it:} \emph{cūḍākaraṇam}
  \item \emph{Say it:} \emph{cūḍākaraṇam}, once more
  \item \emph{Say it:} \emph{cūḍākaraṇe nāpitaḥ śiśoḥ keśān chinatti.} — At the first haircut the barber cuts the baby's hair.
  \item \emph{Recall:} say the Sanskrit for the naming ceremony, then the Sanskrit for a baby's first rice, then say \emph{cūḍākaraṇam} again
\end{itemize}

\begin{tcolorbox}[breakable,colback=teal!4,colframe=teal!35!black,title={Before you move on}]
What does \sk{चूडाकरणम्} mean? (\textquotedblleft{}The first haircut\textquotedblright{}.) Say it once more.
\end{tcolorbox}
\section[\texorpdfstring{\sk{उपनयनम्} (upanayanam)}{उपनयनम् (upanayanam)}]{\sk{उपनयनम्} (upanayanam) — the sacred-thread ceremony}

\label{lesson:SA-C493-upanayanam}



\begin{quote}
Before the new one: say the Sanskrit for a baby's first rice, then the Sanskrit for the first haircut.
\end{quote}

\subsection*{You'll want to know: \sk{उपनयनम्}}
\textbf{\sk{उपनयनम्}} — \emph{upanayanam} — \textquotedblleft{}the sacred-thread ceremony\textquotedblright{}.

\begin{quote}
\sk{उपनयनात्} \sk{अनन्तरं} \sk{बालः} \sk{शिष्यः} \sk{भवति।}
\end{quote}

\begin{quote}
\emph{upanayanāt anantaraṃ bālaḥ śiṣyaḥ bhavati.}
\end{quote}

\begin{quote}
After the sacred-thread ceremony the boy becomes a pupil.
\end{quote}

A neuter noun.

\begin{cousinweb}
\textbf{\sk{उपनयनम्}} (\emph{upanayanam}) is built from \emph{upa}, \textquotedblleft{}near\textquotedblright{}, and the root √nī, \textquotedblleft{}to lead\textquotedblright{}, the root of \textbf{\sk{नयति}}: the teacher leads the boy to himself (MW).
\end{cousinweb}

\begin{grammarlens}[title={What you've built}]
One more word for describing what you lived and felt.
\end{grammarlens}

\subsection*{Your turn}
\begin{itemize}
\raggedright
  \item \emph{Say it:} \emph{upanayanam}
  \item \emph{Say it:} \emph{upanayanam}, once more
  \item \emph{Say it:} \emph{upanayanāt anantaraṃ bālaḥ śiṣyaḥ bhavati.} — After the sacred-thread ceremony the boy becomes a pupil.
  \item \emph{Recall:} say the Sanskrit for a baby's first rice, then the Sanskrit for the first haircut, then say \emph{upanayanam} again
\end{itemize}

\begin{tcolorbox}[breakable,colback=teal!4,colframe=teal!35!black,title={Before you move on}]
What does \sk{उपनयनम्} mean? (\textquotedblleft{}The sacred-thread ceremony\textquotedblright{}.) Say it once more.
\end{tcolorbox}
